\chapter{Rezultati}
\label{sec:rezultati}

\TODO{
\begin{itemize}
    \item To poglavje je začrtano okoli glavnega eksperimenta: binarna klasifikacija valence z $7$-kratnim prečnim preverjanjem po subjektih.
    \item Vir rezultatov je retenirani tek \texttt{RETAIN\_2026-06-12\_16-29-08}.
    \item V besedilu na začetku na kratko zapiši, da ime artefakta še vsebuje staro oznako \texttt{3class}, vendar metapodatki razredov potrjujejo binarno nalogo z razredoma nizka in visoka valenca.
    \item Povej, da so v glavnem besedilu prikazani samo ključni rezultati, razširjene metrike pa so v dodatku~\ref{app:razsirjeni-rezultati-binarna-valenca}.
\end{itemize}
}
% Polna pot vira: GFM-for-eyetracker/results/quick_v1_v2_comparison/RETAIN_2026-06-12_16-29-08.
% Slike rezultatov v tem poglavju so reproducibilno generirane z:
% diploma-latex/skripte/generiraj_slike_rezultatov.py.

\section{Porazdelitev razredov in zanesljivost oznak}
\label{subsec:rezultati-porazdelitev-razredov}

\TODO{
\begin{itemize}
    \item Opiši, da so po odstranitvi nevtralnega razreda razredi nizka/visoka valenca skoraj uravnoteženi pri vseh štirih množicah signalov.
    \item Uporabi številke iz tabele \texttt{label\_distribution\_aggregate.csv}: delež nizke valence je približno $52~\%$, delež visoke valence pa približno $48~\%$.
    \item Poveži večinski klasifikator s porazdelitvijo razredov: visoka točnost večinskega modela ne pomeni dobrega razločevanja, kar pokaže nizek makro F1.
    \item Kratko utemelji, zakaj je binarna valenca glavna ciljna naloga: analiza ujemanja oznak s samoocenami pokaže samo $3{,}3~\%$ neujemanja za binarno valenco, medtem ko je neujemanje pri binarni vzburjenosti približno $29{,}7~\%$.
\end{itemize}
}

% Vir slike: GFM-for-eyetracker/results/quick_v1_v2_comparison/RETAIN_2026-06-12_16-29-08/tables/label_distribution_aggregate.csv.
% Slika je generirana s skriptom diploma-latex/skripte/generiraj_slike_rezultatov.py.
\begin{figure}[H]
    \centering
    \includegraphics[width=0.85\textwidth]{porazdelitev_razredov_delezi_retain_2026-06-12.png}
    \caption{Porazdelitev razredov za glavni binarni eksperiment z valenco.}
    \label{fig:rezultati-porazdelitev-razredov}
\end{figure}

\section{Glavna primerjava modelov}
\label{subsec:rezultati-glavna-primerjava}

\TODO{
\begin{itemize}
    \item Najprej razloži zapis v celicah: prva številka je točnost, druga je makro F1; obe sta v odstotkih.
    \item Poudari, da tabela primerja isto ciljno nalogo za štiri enakovredne množice signalov: samo pogled, samo zenici, pogled in zenici ter vsi signali.
    \item Zapiši glavno ugotovitev: najboljše rezultate dosežejo negrafovski modeli nad agregiranimi značilkami, grafovski modeli pa so večinoma pod najboljšima SVM in LightGBM.
    \item Posebej omeni, da množica pogled in zenici doseže najvišji posamični rezultat, medtem ko dodajanje vseh signalov ne prinese izboljšanja.
    \item Pri \GazeMAE{}/MOMENT pojasni, da je to enoten prikaz zamrznjenih predstavitvenih modelov: \GazeMAE{} pri pogledu, MOMENT pri zenicah in fuzija obeh pri kombiniranih signalih.
\end{itemize}
}

% Samodejno generirano iz GFM-for-eyetracker/results/quick_v1_v2_comparison/RETAIN_2026-06-12_16-29-08/tables/metric_summary_with_std.csv.
% Celice vsebujejo točnost / makro F1 v odstotkih. Standardni odkloni so v dodatku.
\begin{table}[H]
    \centering
    \caption{Glavna primerjava modelov pri binarni klasifikaciji valence. Celice vsebujejo točnost / makro F1 v odstotkih.}
    \label{tab:rezultati-glavna-primerjava}
    \small
    \begin{tabularx}{\textwidth}{@{}>{\raggedright\arraybackslash}p{0.23\textwidth} *{4}{>{\centering\arraybackslash}X}@{}}
        \toprule
        Model & samo pogled & samo zenici & pogled in zenici & vsi signali \\
        \midrule
        Naključni & 50,2 / 50,1 & 50,7 / 50,6 & 50,8 / 50,6 & 50,9 / 50,7 \\
        Večinski & 52,5 / 34,4 & 52,2 / 34,2 & 52,1 / 34,2 & 52,2 / 34,2 \\
        SVM & 67,5 / 67,2 & 62,1 / 60,7 & 70,5 / 69,7 & 68,6 / 67,6 \\
        LightGBM & 69,7 / 69,4 & 63,1 / 61,9 & 69,6 / 68,8 & 69,0 / 68,3 \\
        MLP & 67,0 / 66,7 & 62,0 / 60,3 & 68,7 / 67,6 & 67,2 / 66,2 \\
        GazeMAE/MOMENT & 65,7 / 65,2 & 60,5 / 60,1 & 65,6 / 65,2 & 66,1 / 65,6 \\
        GCN & 64,4 / 64,1 & 61,0 / 58,4 & 65,4 / 63,9 & 59,6 / 57,6 \\
        HeteroGCN-mean & 61,7 / 59,2 & 58,6 / 55,7 & 63,8 / 61,1 & 61,8 / 59,7 \\
        HeteroGCN-MLP & 65,9 / 65,8 & 59,1 / 56,0 & 64,2 / 62,8 & 60,5 / 56,6 \\
        HeteroGCN-MLP-w & 64,3 / 64,0 & 58,4 / 54,5 & 65,5 / 64,4 & 62,7 / 60,8 \\
        \bottomrule
    \end{tabularx}
\end{table}


\TODO{
\begin{itemize}
    \item Spodnja slika prikazuje iste rezultate kot tabela~\ref{tab:rezultati-glavna-primerjava}, vendar v slogu grafov iz eksperimentalnega cevovoda: vrednosti so zapisane kot deleži med $0$ in $1$, barva celice pa sledi točnosti kot glavni metriki.
    \item Po pisanju se odloči, ali v glavnem besedilu obdržiš tabelo, heatmap ali oboje. Če obdržiš oboje, naj besedilo jasno pove, da gre za alternativna prikaza istih rezultatov.
\end{itemize}
}

% Slika je generirana iz:
% GFM-for-eyetracker/results/quick_v1_v2_comparison/RETAIN_2026-06-12_16-29-08/tables/thesis_signal_set_model_metrics.csv.
% Konkretni skript: diploma-latex/skripte/generiraj_slike_rezultatov.py.
\begin{figure}[H]
    \centering
    \includegraphics[width=\textwidth]{glavna_primerjava_modelov_heatmap.png}
    \caption{Toplotni prikaz glavne primerjave modelov. V celicah sta zapisani točnost / makro F1 kot deleža, barva pa prikazuje točnost.}
    \label{fig:rezultati-glavna-primerjava-heatmap}
\end{figure}

\section{Vpliv množice signalov}
\label{subsec:rezultati-signalne-mnozice}

\TODO{
\begin{itemize}
    \item Opiši razliko med štirimi signalnimi množicami in jih obravnavaj enakovredno.
    \item Izpostavi, da je samo zenici šibkejša, vendar informativna signalna množica; ni pomožni ali nedokončan eksperiment.
    \item Primerjaj samo pogled in pogled in zenici: dodajanje zeničnih signalov izboljša najboljši rezultat pri SVM, manj jasno pa pomaga grafovskim modelom.
    \item Pojasni, da vsi signali ne pomenijo nujno boljšega rezultata, saj dodatni signali lahko dodajo šum ali povečajo težavnost modeliranja.
    \item Pri interpretaciji ne sklepaj, da so zenice ali razdalja nepomembne na splošno; sklep velja za uporabljeno predstavitev, podatkovno zbirko in protokol.
\end{itemize}
}

\section{Matrike zmede}
\label{subsec:rezultati-matrike-zmede}

\TODO{
\begin{itemize}
    \item Matrike zmede so prikazane samo za množico pogled in zenici, ker tam nastopi najboljši skupni rezultat.
    \item Primerjaj najboljši splošni model SVM in najboljši grafovski model, tj. končni GNN.
    \item Opiši, kateri razred je pogosteje zamenjan in ali model kaže izrazit zamik proti nizki ali visoki valenci.
    \item Povej, da so prikazane samo vrstično normalizirane vrednosti; absolutne matrike zmede sodijo v dodatek, če jih bomo sploh vključili.
\end{itemize}
}

% Matrike so generirane iz:
% GFM-for-eyetracker/results/quick_v1_v2_comparison/RETAIN_2026-06-12_16-29-08/tables/confusion_matrices.csv.
% Konkretni skript: diploma-latex/skripte/generiraj_slike_rezultatov.py.
\begin{figure}[H]
    \centering
    \includegraphics[width=\textwidth]{matrike_zmede_pogled_zenici_svm_heterogcn_mlp_w.png}
    \caption{Vrstično normalizirani matriki zmede za najboljši model SVM in najboljši grafovski model HeteroGCN-MLP-w pri množici signalov pogled in zenici.}
    \label{fig:rezultati-matrike-zmede}
\end{figure}

\section{Velikost in računska zahtevnost modelov}
\label{subsec:rezultati-racunska-zahtevnost}

\TODO{
\begin{itemize}
    \item Tabela naj bo uporabljena za praktično interpretacijo rezultatov: ali bolj zapleteni modeli prinesejo dovolj boljše rezultate glede na svojo ceno.
    \item Pojasni, da za SVM, LightGBM, večinski in naključni model število nevronskih parametrov ni smiselna primerjava, zato je označeno z \texttt{--}.
    \item Pri \GazeMAE{}/MOMENT loči učljive parametre glave in vse parametre zamrznjenih kodirnikov.
    \item Izpostavi, da so grafovski modeli računsko dražji na okno kot tabelarični modeli, čeprav v glavnih metrikah ne dosežejo najboljših rezultatov.
\end{itemize}
}

% Samodejno generirano iz .../tables/main_model_complexity_report.csv za signalno množico pogled in zenici.
\begin{table}[H]
    \centering
    \caption{Velikost in računska zahtevnost modelov pri množici signalov pogled in zenici.}
    \label{tab:rezultati-racunska-zahtevnost}
    \scriptsize
    \setlength{\tabcolsep}{2pt}
    \begin{tabular}{@{}
    >{\raggedright\arraybackslash}p{0.22\textwidth}
    >{\raggedright\arraybackslash}p{0.12\textwidth}
    >{\raggedleft\arraybackslash}p{0.08\textwidth}
    >{\raggedleft\arraybackslash}p{0.10\textwidth}
    >{\raggedleft\arraybackslash}p{0.10\textwidth}
    >{\raggedleft\arraybackslash}p{0.11\textwidth}
    >{\raggedleft\arraybackslash}p{0.135\textwidth}
    @{}}
        \toprule
        Model &
        Vhod &
        \makecell[r]{Učljivi\\{[k]}} &
        \makecell[r]{Vsi\\{[k]}} &
        \makecell[r]{Učenje\\{[ms/ok.]}} &
        \makecell[r]{Inferenca\\{[ms/ok.]}} &
        \makecell[r]{Točnost/\\makro F1} \\
        \midrule
        Naključni & brez značilk & -- & -- & 0,00 & 0,00 & 50,8/50,6 \\
        Večinski & brez značilk & -- & -- & 0,00 & 0,00 & 52,1/34,2 \\
        \midrule
        \ModelName{SVM} & agreg. znač. & -- & -- & 0,39 & 0,08 & \textbf{70,5/69,7} \\
        \ModelName{LightGBM} & agreg. znač. & -- & -- & 0,04 & 0,00 & 69,6/68,8 \\
        \ModelName{MLP} & agreg. znač. & 7 & 7 & 0,39 & 0,00 & 68,7/67,6 \\
        \midrule
        \ModelName{GazeMAE/MOMENT} & signal & 181 & 127.217 & 0,26 & 0,00 & 65,6/65,2 \\
        \midrule
        \ModelName{GCN} & grafi & 26 & 26 & 20,73 & 0,19 & 65,4/63,9  \\
        \ModelName{HeteroGCN-mean} & grafi & 43 & 43 & 26,64 & 0,20 & 63,8/61,1  \\
        \ModelName{HeteroGCN-MLP} & grafi & 76 & 76 & 22,27 & 0,20 & 64,2/62,8 \\
        \ModelName{HeteroGCN-MLP-w} & grafi & 76 & 76 & 21,74 & 0,19 & \textbf{\emph{65,5/64,4}} \\
        \bottomrule
    \end{tabular}
\end{table}


\section{Dodatna preverjanja}
\label{subsec:rezultati-dodatna-preverjanja}

\TODO{
\begin{itemize}
    \item Kratko omeni dodatni poskus izbire grafovske konvolucije iz dodatka~\ref{app:izbira-grafovske-konvolucije}.
    \item Povej samo bistvo: \texttt{GCNConv} ni bil absolutno najboljši operator v dodatnem poskusu, vendar je bil dovolj stabilen in poenostavi primerjavo arhitekturnih stopenj.
    \item Ne razširjaj tega dela v novo glavno primerjavo modelov; podrobne številke naj ostanejo v dodatku.
\end{itemize}
}

\section{Povzetek rezultatov}
\label{subsec:rezultati-povzetek}

\TODO{
\begin{itemize}
    \item V nekaj odstavkih povzemi odgovore na raziskovalno vprašanje.
    \item Prvi odstavek: najboljši rezultati pripadajo negrafovskim modelom, predvsem SVM in LightGBM.
    \item Drugi odstavek: grafovska predstavitev dosega rezultate nad naključnim in večinskim izhodiščem, vendar trenutna grafovska arhitektura ne izboljša klasifikacije glede na agregirane tabelarične značilke.
    \item Tretji odstavek: signalna množica pogled in zenici je najboljša v tej evalvaciji, samo zenici pa ostane smiselna samostojna primerjava.
    \item Zaključi previdno: rezultati ne zavračajo uporabnosti grafovskega modeliranja podatkov sledilnika pogleda, kažejo pa, da uporabljena grafovska arhitektura in konstrukcija povezav v tej nalogi nista presegli močnejših negrafovskih primerjav.
\end{itemize}
}
